\documentclass[../DoAn.tex]{subfiles}
\begin{document}

\section{Khảo sát các giải pháp ẩn danh hiện tại}
Để xác định ưu điểm, hạn chế và khoảng trống công nghệ của các giải pháp bảo mật trên nền tảng Ethereum, đồ án tiến hành khảo sát ba hướng tiếp cận tiêu biểu đại diện cho các mô hình ẩn danh khác nhau.

\subsection{Máy trộn tiền phi tập trung: Tornado Cash}

Tornado Cash là một giao thức bảo vệ quyền riêng tư trên Ethereum sử dụng công nghệ bằng chứng không tiết lộ thông tin (Zero-Knowledge Proof - ZK-SNARK) nhằm phá vỡ mối liên hệ trực tiếp giữa địa chỉ gửi và địa chỉ nhận. Người dùng gửi tài sản vào một Pool chung của hợp đồng thông minh và nhận về một mã bí mật (secret note). Khi muốn rút tiền, người dùng tạo bằng chứng ZK để chứng minh quyền sở hữu một khoản tiền hợp lệ trong Pool mà không tiết lộ khoản tiền đó được gửi từ giao dịch nào~\cite{tornado_docs}.

\begin{itemize}
\item \textbf{Ưu điểm:} Cung cấp mức độ riêng tư cao nhờ cơ chế ZK-SNARK và tập ẩn danh (Anonymity Set) lớn. Khi số lượng người tham gia trong Pool tăng lên, việc liên kết giữa giao dịch nạp và giao dịch rút trở nên cực kỳ khó khăn~\cite{tornado_docs}.


\item \textbf{Hạn chế:} Hiệu quả bảo mật phụ thuộc mạnh vào kích thước của tập ẩn danh. Ngoài ra, giao thức chủ yếu hỗ trợ các giao dịch nạp/rút tài sản và chưa hỗ trợ trực tiếp các hoạt động DeFi riêng tư phức tạp. Về mặt pháp lý, Tornado Cash cũng đối mặt với nhiều tranh cãi do có thể bị lợi dụng cho các hoạt động rửa tiền và che giấu nguồn gốc tài sản bất hợp pháp~\cite{tornado_docs}.


\end{itemize}

\subsection{Giao thức địa chỉ ẩn danh: Umbra Cash}

Umbra Cash là một giao thức hiện thực hóa cơ chế Stealth Address trên Ethereum. Người nhận công khai một Meta-Address, từ đó người gửi có thể tạo ra một địa chỉ nhận duy nhất (one-time stealth address) cho mỗi giao dịch. Chỉ người nhận mới có khả năng xác định và kiểm soát địa chỉ này, giúp che giấu danh tính người nhận trên chuỗi khối mà không cần sử dụng Pool trộn tài sản~\cite{umbra_docs,erc5564}.

\begin{itemize}
\item \textbf{Ưu điểm:} Cho phép chuyển tài sản trực tiếp giữa người gửi và người nhận mà không cần thông qua Pool trung gian, giúp giảm chi phí giao dịch và phù hợp với các kịch bản thanh toán riêng tư như chuyển lương hoặc thanh toán P2P~\cite{umbra_docs}.


\item \textbf{Hạn chế:} Mức độ riêng tư phụ thuộc vào số lượng người dùng sử dụng giao thức và có thể bị suy giảm bởi các kỹ thuật phân tích hành vi giao dịch (Deanonymization). Mặc dù Umbra đã sử dụng cơ chế Relayer để hỗ trợ chi trả phí gas cho một số trường hợp, giao thức vẫn phụ thuộc vào hạ tầng relayer và chưa tận dụng đầy đủ các khả năng của Account Abstraction như Paymaster, Social Recovery hay Session Keys~\cite{umbra_docs}.

\end{itemize}

\subsection{Giao thức tài khoản riêng tư: Nocturne}

Nocturne là một giao thức quyền riêng tư trên Ethereum kết hợp giữa bằng chứng không tiết lộ thông tin (ZKP) và cơ chế Account Abstraction theo tiêu chuẩn ERC-4337~\cite{erc4337}. Thay vì chỉ ẩn danh cho từng giao dịch riêng lẻ, Nocturne xây dựng mô hình tài khoản riêng tư (Private Account), cho phép người dùng thực hiện nhiều hoạt động như lưu trữ, chuyển đổi và quản lý tài sản trong một môi trường bảo mật thống nhất~\cite{nocturne_docs}.

\begin{itemize}
\item \textbf{Ưu điểm:} Kết hợp được quyền riêng tư của ZKP với khả năng cải thiện trải nghiệm người dùng từ Account Abstraction, bao gồm hỗ trợ thanh toán phí gas linh hoạt, tài trợ phí giao dịch thông qua Paymaster và khả năng mở rộng cho các ứng dụng DeFi riêng tư~\cite{nocturne_docs}.


\item \textbf{Hạn chế:} Kiến trúc hệ thống phức tạp hơn đáng kể so với các giải pháp truyền thống do phải phối hợp giữa các thành phần mật mã học, Account Abstraction và cơ chế quản lý trạng thái riêng tư. Điều này làm tăng chi phí tính toán và độ phức tạp trong triển khai hệ thống~\cite{nocturne_docs}.


\end{itemize}

\textbf{Kết luận khảo sát:} Các giải pháp hiện nay tiếp cận bài toán quyền riêng tư theo những hướng khác nhau. Tornado Cash cung cấp mức độ ẩn danh cao nhưng gặp các vấn đề về khả năng sử dụng và tuân thủ. Umbra Cash đơn giản hơn và phù hợp cho các giao dịch riêng tư trực tiếp nhưng vẫn tồn tại nguy cơ deanonymization và phụ thuộc vào relayer. Nocturne mở rộng quyền riêng tư lên cấp độ tài khoản và tận dụng Account Abstraction để cải thiện trải nghiệm người dùng, tuy nhiên hệ thống có độ phức tạp triển khai cao. Từ đó, đồ án định hướng xây dựng một kiến trúc kết hợp giữa Stealth Address, Account Abstraction và ZK-SNARK nhằm tận dụng ưu điểm của từng mô hình, đồng thời khắc phục các hạn chế về khả năng sử dụng, chi trả phí gas và bảo vệ quyền riêng tư trên Ethereum~\cite{erc5564,erc4337,groth16}.

\section{Phân tích yêu cầu hệ thống}

Dựa trên những định hướng từ phần khảo sát, đồ án tiến hành đặc tả các yêu cầu cốt lõi nhằm định hình kiến trúc phần mềm.

\subsection{Yêu cầu chức năng}

Hệ thống cần cung cấp đầy đủ các chức năng phục vụ việc quản lý danh tính, thực hiện giao dịch riêng tư và bảo vệ tài khoản của người dùng trên nền tảng blockchain.

\begin{enumerate}

\item \textbf{Quản lý danh tính người dùng}
\begin{itemize}
\item \textbf{FR-01:} Cho phép người dùng tạo mới danh tính và ví ẩn danh.
\item \textbf{FR-02:} Cho phép lưu trữ và quản lý thông tin định danh phục vụ cho các giao dịch riêng tư.
\item \textbf{FR-03:} Cho phép khôi phục hoặc nhập lại ví từ thông tin xác thực đã được lưu trước đó.
\item \textbf{FR-04:} Hỗ trợ liên kết danh tính người dùng với các dịch vụ định danh phi tập trung.
\end{itemize}

\item \textbf{Thực hiện giao dịch riêng tư}
\begin{itemize}
\item \textbf{FR-05:} Cho phép người dùng gửi tài sản đến người nhận thông qua định danh công khai hoặc định danh phi tập trung.
\item \textbf{FR-06:} Hệ thống phải tạo ra các địa chỉ nhận riêng biệt cho từng giao dịch nhằm giảm khả năng liên kết giữa các giao dịch.
\item \textbf{FR-07:} Hỗ trợ chuyển nhiều loại tài sản số khác nhau như tiền mã hóa, token và tài sản số đại diện.
\end{itemize}

\item \textbf{Nhận diện và quản lý tài sản}
\begin{itemize}
\item \textbf{FR-08:} Cho phép người dùng phát hiện các giao dịch và tài sản thuộc quyền sở hữu của mình.
\item \textbf{FR-09:} Hệ thống phải xác thực quyền sở hữu đối với các tài sản được phát hiện.
\item \textbf{FR-10:} Hiển thị thông tin tài sản và số dư thuộc quyền quản lý của người dùng.
\end{itemize}

\item \textbf{Chi tiêu tài sản riêng tư}
\begin{itemize}
\item \textbf{FR-11:} Cho phép người dùng thực hiện các giao dịch chuyển tài sản từ tài khoản riêng tư.
\item \textbf{FR-12:} Hệ thống phải xác minh quyền chi tiêu của người dùng mà không làm lộ các thông tin nhạy cảm.
\item \textbf{FR-13:} Hỗ trợ cơ chế thanh toán phí giao dịch linh hoạt nhằm nâng cao trải nghiệm sử dụng.
\item \textbf{FR-14:} Đảm bảo người dùng có thể sử dụng tài khoản riêng tư mà không cần thực hiện các bước cấu hình phức tạp.
\end{itemize}

\item \textbf{Phục hồi tài khoản}
\begin{itemize}
\item \textbf{FR-15:} Cho phép người dùng thiết lập cơ chế phục hồi tài khoản thông qua các bên được ủy quyền.
\item \textbf{FR-16:} Cho phép đề xuất cập nhật thông tin xác thực khi xảy ra sự cố mất quyền truy cập.
\item \textbf{FR-17:} Hệ thống phải hỗ trợ quy trình xác nhận và thực hiện phục hồi tài khoản theo các điều kiện đã thiết lập.
\item \textbf{FR-18:} Đảm bảo dữ liệu phục hồi được đồng bộ và nhất quán giữa các thành phần của hệ thống.
\end{itemize}

\end{enumerate}

\subsection{Yêu cầu phi chức năng}

\begin{itemize}

\item \textbf{Tính riêng tư (Privacy):}
Hệ thống phải giảm thiểu khả năng liên kết giữa người gửi, người nhận và các giao dịch trên chuỗi khối nhằm bảo vệ danh tính và lịch sử giao dịch của người dùng.

\item \textbf{Bảo mật (Security):}
Hệ thống phải bảo vệ các thông tin nhạy cảm và đảm bảo chỉ chủ sở hữu hợp lệ mới có quyền truy cập hoặc thực hiện các thao tác đối với tài sản của mình.

\item \textbf{Tính đúng đắn (Correctness):}
Mọi giao dịch và yêu cầu thực hiện trên hệ thống phải được xác thực hợp lệ trước khi được chấp nhận xử lý.

\item \textbf{Hiệu năng (Performance):}
Hệ thống cần đáp ứng thời gian xử lý phù hợp cho các tác vụ chính nhằm đảm bảo trải nghiệm người dùng.

\item \textbf{Khả năng mở rộng (Scalability):}
Hệ thống cần hỗ trợ số lượng người dùng và giao dịch tăng lên mà không làm suy giảm đáng kể hiệu năng hoạt động.

\item \textbf{Độ tin cậy (Reliability):}
Hệ thống phải duy trì tính ổn định, hạn chế mất mát dữ liệu và đảm bảo tính nhất quán giữa các thành phần.

\item \textbf{Khả năng sử dụng (Usability):}
Giao diện và quy trình sử dụng cần đơn giản, dễ tiếp cận đối với người dùng không có kiến thức chuyên sâu về blockchain.

\end{itemize}

\section{Biểu đồ Use Case tổng quát}
Hệ thống được thiết kế xoay quanh sự tương tác của ba tác nhân chính:
\begin{itemize}
    \item \textbf{Người nhận (Receiver/Owner):} Thiết lập ví ẩn danh, nhận tài sản và thực hiện các thao tác chi tiêu hoặc phục hồi quyền truy cập.
    
    \item \textbf{Người gửi (Sender):} Tính toán địa chỉ Stealth và gửi tài sản đến người nhận.
    
    \item \textbf{Người giám hộ (Guardian):} Tham gia xác nhận quá trình phục hồi quyền truy cập của người dùng.
\end{itemize}

\begin{figure}[H]
    \centering
    \includegraphics[width=0.6\textwidth]{SOICT_DATN_NGUYEN_VAN_AN_20224910 (1)/Hinhve/usecase_diagram_123.png}
    \caption{Biểu đồ Use Case tổng quát của hệ thống Stealth Address}
    \label{fig:UseCase}
\end{figure}

\end{document}